\section{The distribution under logarithmic averaging}\label{sec:distribution}
Put $b_j=1+\log a_j$ for $j=p,q$,
$B=b_p+b_q$, and $A=a_p+a_q$. Define the two positive constants
\begin{equation}
 h_j=\frac{1}{\ell_j}\left(\frac{1}{1-a_j}+
                         \frac{1}{j a_j-1}\right),\qquad j=p,q .
 \label{eq:h}
\end{equation}
\begin{theorem}\label{thm:log}
Let $U,V$ be independent uniform variables on $\T$ and let
$Z=f_p(U)+f_q(V)$. For every continuous $\varphi$ on $[B,A]$,
\begin{equation}
 \lim_{X\to\infty}\frac1{\log X}\int_1^X
   \varphi(R_{p,q}(x))\frac{\mathrm dx}{x}
   =\E\varphi(Z). \label{eq:loglaw}
\end{equation}
The support is $[B,A]$, the law is absolutely continuous, and
$\E Z=2-\tfrac12\log(pq)$. Its tails satisfy, as $s\downarrow0$,
\begin{equation}
 \Prob(Z\leq B+s)\sim \frac{2\pi}{\ell_p\ell_q}s,\qquad
 \Prob(Z\geq A-s)\sim \frac{h_p h_q}{2}s^2.
 \label{eq:tails}
\end{equation}
Furthermore $\liminf_{x\to\infty}R_{p,q}(x)=B$ and
$\limsup_{x\to\infty}R_{p,q}(x)=A$.
\end{theorem}
\begin{proof}
Unique factorization gives $\ell_p/\ell_q\notin\mathbb Q$.
Thus $k\cdot\omega\ne0$ for every nonzero $k\in\mathbb Z^2$.
With $t=\log x$, the time average on $0\leq t\leq T$ of the character
$e^{2\pi i (k\cdot\omega)t}$ tends to zero. Uniform trigonometric
approximation of $\varphi\circ F$ proves \eqref{eq:loglaw}. This is
the continuous-time character argument for equidistribution \cite{weyl}.

For $j=p,q$ the inequalities $(j-1)/j<\ell_j<j-1$ give
$a_j<1<j a_j$. The unique minimum of $f_j$ is $b_j$ at
$v_j^*=-\log a_j/\ell_j$ and its maximum is $a_j$ at the joined
endpoints. The map $F$ attains every value in $[B,A]$, and the inverse
image of any relatively open subinterval has positive area. The same
equidistribution argument then gives the stated extremal limits.
Direct integration gives $\int_0^1 f_j(v)\dd v=1-\ell_j/2$.

For $b_j<y<a_j$, the density of $f_j(U)$ is
\begin{equation}
 g_j(y)=\sum_{v\in(0,1):\, f_j(v)=y}
               \frac{1}{\ell_j\,|a_j e^{\ell_j v}-1|}.
 \label{eq:density}
\end{equation}
This follows by changing variables on the two strictly monotone
branches. The critical point and the endpoints have zero measure.
Independence gives the convolution $g_p*g_q$ for the density of $Z$.

At the minimum, $f_j(v_j^*+w)=b_j+\ell_j^2w^2/2+O(w^3)$.
Consequently $g_j(b_j+u)\sim\sqrt{2}/(\ell_j\sqrt{u})$. At the
maximum the two one-sided derivatives are nonzero, giving
$g_j(a_j-u)\to h_j$. For the lower tail, integrate the product of
the two shifted densities over the triangle $u,v>0$, $u+v\leq s$.
After rescaling by $s$, the singular factors have integral
\[
 \int_0^1\int_0^{1-u}(uv)^{-1/2}\dd v\dd u=\pi .
\]
The local asymptotics bound the rescaled integrand by a constant
multiple of $(uv)^{-1/2}$. Dominated convergence yields the first
constant in \eqref{eq:tails}. At the upper end the two densities are
bounded and the triangle has area $s^2/2$, giving the second.
\end{proof}

For $m=24$ the endpoints and mean become
\[
 B=2+\log\!\left(\frac{\log 2\,\log 3}{2}\right),\qquad
 A=\log 2+\tfrac12\log 3,\qquad \E Z=2-\tfrac12\log 6.
\]
Their values are approximately $1.034388$, $1.242453$ and $1.104120$.
The square-root singularity of each local density at its minimum
accounts for the linear lower tail. At the maximum both slopes of each
local function are nonzero, producing the quadratic upper tail.
